%-------------------------
% Ashir Borah — Academic CV
% Adapted from Jake Gutierrez's CV template (MIT)
%   https://github.com/jakegut/resume
% Build: `make` (uses pdflatex; see Makefile)
%-------------------------
\documentclass[letterpaper,11pt]{article}

\usepackage{latexsym}
\usepackage[empty]{fullpage}
\usepackage{titlesec}
\usepackage{marvosym}
\usepackage[usenames,dvipsnames]{color}
\usepackage{verbatim}
\usepackage{enumitem}
\usepackage[hidelinks]{hyperref}
\usepackage{fancyhdr}
\usepackage[english]{babel}
\usepackage{tabularx}
\usepackage{xcolor}
\usepackage{multicol}
\input{glyphtounicode}

\pagestyle{fancy}
\fancyhf{}
\fancyfoot[L]{\small\itshape Ashir Borah — CV}
\fancyfoot[R]{\small\thepage}
\renewcommand{\headrulewidth}{0pt}
\renewcommand{\footrulewidth}{0pt}
\setlength{\footskip}{14pt}

% margins
\addtolength{\oddsidemargin}{-0.5in}
\addtolength{\evensidemargin}{-0.5in}
\addtolength{\textwidth}{1in}
\addtolength{\topmargin}{-0.5in}
\addtolength{\textheight}{1.0in}

\urlstyle{same}

\raggedbottom
\raggedright
\setlength{\tabcolsep}{0in}

% section formatting
\titleformat{\section}{
  \vspace{-4pt}\scshape\raggedright\large
}{}{0em}{}[\color{black}\titlerule \vspace{-5pt}]

% machine-readable PDF
\pdfgentounicode=1

%-------------------------
% Custom commands
%-------------------------
\newcommand{\me}{\textbf{A.A.~Borah}}
\newcommand{\etal}{\textit{et~al}}

\newcommand{\resumeItem}[1]{%
  \item\small{#1 \vspace{-2pt}}
}

\newcommand{\cvHeading}[4]{%
  \vspace{-2pt}\item
    \begin{tabular*}{0.97\textwidth}[t]{l@{\extracolsep{\fill}}r}
      \textbf{#1} & #2 \\
      \textit{\small#3} & \textit{\small #4} \\
    \end{tabular*}\vspace{-7pt}
}

\newcommand{\cvSubheading}[2]{%
  \item
    \begin{tabular*}{0.97\textwidth}{l@{\extracolsep{\fill}}r}
      \small#1 & \textit{\small #2} \\
    \end{tabular*}\vspace{-7pt}
}

\newcommand{\cvProject}[2]{%
  \item
    \begin{tabular*}{0.97\textwidth}{l@{\extracolsep{\fill}}r}
      \small\textbf{#1} & \textit{\small #2} \\
    \end{tabular*}\vspace{-7pt}
}

% publication entry: {authors}{title}{venue}{year}{doi}
% DOI is the bare identifier (e.g., 10.1038/s41586-020-2769-8).
\newcommand{\pub}[5]{%
  \item\small{
    #1.\ \textbf{#2}.\ \textit{#3} (#4).\ \href{https://doi.org/#5}{doi:#5}
    \vspace{-2pt}
  }
}

% sub-grouping label inside the Publications section
\newcommand{\pubGroup}[1]{%
  \vspace{4pt}\noindent\textit{\textbf{#1}}\vspace{-4pt}
}

% conference abstract / poster: {authors}{title}{venue with vol/issue/abstract id}{year}
% No DOI required; AACR & similar abstracts are usually cited by abstract ID alone.
\newcommand{\abstr}[4]{%
  \item\small{
    #1.\ \textbf{#2}.\ \textit{#3} (#4).
    \vspace{-2pt}
  }
}

\renewcommand\labelitemii{$\vcenter{\hbox{\tiny$\bullet$}}$}

\newcommand{\sectionList}{\begin{itemize}[leftmargin=0.15in, label={}]}
\newcommand{\sectionListEnd}{\end{itemize}}
\newcommand{\bullets}{\begin{itemize}}
\newcommand{\bulletsEnd}{\end{itemize}\vspace{-5pt}}

%======================================
\begin{document}

%-------------------------
% Header
%-------------------------
\begin{center}
    {\Huge \scshape Ashir A. Borah} \\ \vspace{2pt}
    \small
    San Francisco, CA \;|\;
    \href{mailto:ashiraseesh@gmail.com}{ashiraseesh@gmail.com} \;|\;
    +1 717-440-0204 \\
    \href{https://ashirborah.com}{ashirborah.com} \;|\;
    \href{https://scholar.google.com/citations?user=HztaA54AAAAJ}{Google Scholar} \;|\;
    \href{https://orcid.org/0000-0003-1193-4987}{ORCID 0000-0003-1193-4987} \;|\;
    \href{https://www.linkedin.com/in/ashirborah}{LinkedIn} \;|\;
    \href{https://github.com/AshirBorah}{GitHub}\\
    \vspace{2pt}
    \textit{\footnotesize Updated April 2026}
\end{center}

%-------------------------
% Education
%-------------------------
\section{Education}
\sectionList

  \cvHeading
    {University of California, San Francisco}{Sep 2022 -- Present}
    {Ph.D.\ Candidate, Biological and Medical Informatics}{San Francisco, CA}
    \cvSubheading{\textbf{M.S., Biological and Medical Informatics} --- awarded en route to Ph.D.\ (advisor: Luke Gilbert)}{Jun 2024}
    \bullets
      \resumeItem{Advisors: Luke Gilbert (UCSF / Arc Institute) and Brian Hie (Stanford / Arc Institute).}
      \resumeItem{Research at the intersection of functional genomics, virology, and computational biology.}
    \bulletsEnd

  \cvHeading
    {Dickinson College}{Aug 2015 -- May 2019}
    {Bachelor of Science, Mathematics and Computer Science}{Carlisle, PA}
    \bullets
      \resumeItem{\textit{Magna Cum Laude}, Phi Beta Kappa. GPA: 3.89/4.00 (Major: 3.90/4.00).}
      \resumeItem{Senior research: \textit{Computational analysis of phorbol 12-myristate 13-acetate (PMA) perturbation on human leukemia} (advisor: Michael Roberts).}
    \bulletsEnd

\sectionListEnd

%-------------------------
% Research Experience
%-------------------------
\section{Research Experience}
\sectionList

  \cvHeading
    {Arc Institute}{Jun 2023 -- Present}
    {Ph.D.\ Researcher (joint with UCSF)}{Palo Alto, CA}
    \bullets
      \resumeItem{Design and develop experimental and computational methods for functional genomics and virology.}
      \resumeItem{Build CRISPR-based screening and perturbation systems to study gene regulation, host--virus biology, and disease mechanisms.}
      \resumeItem{Develop assay platforms and analysis workflows for high-dimensional genomic datasets, including Perturb-seq.}
      \resumeItem{Build AI-enabled scientific tooling --- LLM-based councils, research agents, APIs, and multi-agent workflows for literature synthesis, hypothesis generation, and experimental planning.}
    \bulletsEnd

  \cvHeading
    {Broad Institute of MIT and Harvard}{Jul 2019 -- Apr 2022}
    {Computational Associate II (2021--2022); Computational Associate I (2019--2021)}{Cambridge, MA}
    \cvSubheading{Cancer Data Science --- Cancer Dependency Map (DepMap) Project --- advisor: James McFarland}{}
    \bullets
      \resumeItem{Led the computational effort to validate therapeutic genetic targets and identify their biomarkers.}
      \resumeItem{Developed strategies for discovering new targets from genome-scale CRISPR knockout screens, with a focus on GI cancers.}
      \resumeItem{Contributed to an open-source machine learning platform that models genetic-knockout dependency profiles from $>$100{,}000 genomic features.}
      \resumeItem{Modeled latent variables from genome-scale CRISPR-Cas9 screens to discover a new Integrator complex interaction.}
      \resumeItem{Analyzed RNAi and CRISPR screens as complementary data modalities for cancer-vulnerability discovery.}
    \bulletsEnd

  \cvHeading
    {Dickinson College}{2018 -- 2019}
    {Computational Research Assistant, Biochemistry and Molecular Biology}{Carlisle, PA}
    \cvSubheading{Advisor: Michael Roberts --- topic: PMA perturbation in human leukemia}{}
    \bullets
      \resumeItem{Processed CRISPR knockout and overexpression genomic data on HL-60 leukemia cell lines.}
      \resumeItem{Identified differentially expressed genes from a compound perturbation that induced differentiation.}
      \resumeItem{Presented findings at the 2020 American Association for Cancer Research (AACR) annual meeting; subsequent students followed up on additional targets.}
    \bulletsEnd

  \cvHeading
    {Uliza}{Sep 2017 -- Dec 2017}
    {Machine Learning Developer}{Remote}
    \bullets
      \resumeItem{Built ML-driven systems to support crowdsourced workflow automation.}
      \resumeItem{Helped set up core technical infrastructure --- server management, database optimization, and disaster recovery.}
    \bulletsEnd

\sectionListEnd

%-------------------------
% Publications
%-------------------------
\section{Publications}
\vspace{2pt}
\small
\textbf{Citation metrics (Google Scholar, Apr 2026):} 1{,}125 citations \;|\; h-index 10 \;|\; i10-index 10. \\
Author names with \me\ are mine.
\normalsize
\vspace{2pt}

% Publications list — included from cv.tex.
% Newest first within each subsection.
% \pub args:   {authors}{title}{venue}{year}{doi}      (bare DOI, no URL prefix)
% \abstr args: {authors}{title}{venue+id}{year}        (no DOI; for conference abstracts)
%
% TODO: refresh citation metrics in cv.tex when Scholar updates.

\pubGroup{Peer-reviewed publications}
\sectionList

  \pub
    {T.A. O'Loughlin, J.S. Stiles, P. Acharya, A. Arab, L. Goudy, R. Dai, \me, \etal}
    {mTORC1 activity suppresses ferroptosis through a SCARB1-dependent HDL-tocopherol uptake pathway}
    {Molecular Cell}
    {2026}
    {10.1016/j.molcel.2026.02.019}

  \pub
    {H. Karner, T.C. Mittmann, V.W. Chen, \me, A. Langen, H. Yousefi, L. Fish, \etal}
    {Integrative analysis of mRNA stability regulation uncovers a metastasis-suppressive program in breast cancer}
    {Science Advances}
    {2026}
    {10.1126/sciadv.aea9061}

  \pub
    {J. Gencel-Augusto, H. Li, L.C. Woerner, N. Tian, \me, J.N. Myers, P. Ha, \etal}
    {Human Papillomavirus does not fully inactivate p53 cellular activity in HNSCC}
    {Head \& Neck}
    {2026}
    {10.1002/hed.70085}

  \pub
    {P.C. Borck, I. Boyle, K. Jankovic, N. Bick, K. Foster, A.C. Lau, L.I. Parker-Burns, \ldots, \me, \etal}
    {SKI complex loss renders 9p21.3-deleted or MSI-H cancers dependent on PELO}
    {Nature}
    {2025}
    {10.1038/s41586-024-08509-3}

  \pub
    {L. Goudy, A. Ha, \me, J.M. Umhoefer, L. Chow, C. Tran, A. Winters, \etal}
    {Integrated epigenetic and genetic programming of primary human T cells}
    {Nature Biotechnology}
    {2025}
    {10.1038/s41587-025-02856-w}

  \pub
    {M. Khoroshkin, D. Asarnow, S. Zhou, A. Navickas, A. Winters, J. Goudreau, \ldots, \me, \etal}
    {A systematic search for RNA structural switches across the human transcriptome}
    {Nature Methods}
    {2024}
    {10.1038/s41592-024-02335-1}

  \pub
    {S.J. Liu, C. Zou, J. Pak, A. Morse, D. Pang, T. Casey-Clyde, \me, D. Wu, \etal}
    {In vivo Perturb-seq of cancer and microenvironment cells dissects oncologic drivers and radiotherapy responses in glioblastoma}
    {Genome Biology}
    {2024}
    {10.1186/s13059-024-03404-6}

  \pub
    {V.B. Venkadakrishnan, A.G. Presser, R. Singh, M.A. Booker, N.A. Traphagen, \ldots, \me, \etal}
    {Lineage-specific canonical and non-canonical activity of EZH2 in advanced prostate cancer subtypes}
    {Nature Communications}
    {2024}
    {10.1038/s41467-024-51156-5}

  \pub
    {L.D. Cervia, T. Shibue, \me, B. Gaeta, L. He, L. Leung, N. Li, \etal}
    {A ubiquitination cascade regulating the integrated stress response and survival in carcinomas}
    {Cancer Discovery}
    {2023}
    {10.1158/2159-8290.CD-22-1230}

  \pub
    {J.M. Krill-Burger, J.M. Dempster, \me, B.R. Paolella, D.E. Root, T.R. Golub, \etal}
    {Partial gene suppression improves identification of cancer vulnerabilities when CRISPR-Cas9 knockout is pan-lethal}
    {Genome Biology}
    {2023}
    {10.1186/s13059-023-03020-w}

  \pub
    {J. Pan, J.J. Kwon, J.A. Talamas, \me, F. Vazquez, J.S. Boehm, \etal}
    {Sparse dictionary learning recovers pleiotropy from human cell fitness screens}
    {Cell Systems}
    {2022}
    {10.1016/j.cels.2022.01.005}

  \pub
    {S. Raghavan, P.S. Winter, A.W. Navia, H.L. Williams, A. DenAdel, K.E. Lowder, \ldots, \me, \etal}
    {Microenvironment drives cell state, plasticity, and drug response in pancreatic cancer}
    {Cell}
    {2021}
    {10.1016/j.cell.2021.11.017}

  \pub
    {N. van Wietmarschen, S. Sridharan, W.J. Nathan, A. Tubbs, E.M. Chan, \ldots, \me, \etal}
    {Repeat expansions confer WRN dependence in microsatellite-unstable cancers}
    {Nature}
    {2020}
    {10.1038/s41586-020-2769-8}

\sectionListEnd

\pubGroup{Preprints}
\sectionList

  \pub
    {B.J. Woo, S. Sobti, J.M. Suh, H. Yousefi, K. Garcia, S. Zhou, \me, \etal}
    {Systematic identification of chromatin organizers as tuners of intratumoral heterogeneity}
    {bioRxiv}
    {2026}
    {10.1101/2026.04.18.719392}

\sectionListEnd

\pubGroup{Conference abstracts \& posters}
\sectionList

  \abstr
    {P.C. Borck, I. Boyle, K. Jankovic, N. Bick, K. Foster, A. Lau, L. Parker-Burns, \ldots, \me, \etal}
    {Exploiting dysregulated ribosomal homeostasis in chromosome 9p21.3-deleted cancers and microsatellite-unstable cancers}
    {Cancer Research 85 (8 Suppl.\ 2), Abstract SY14-03; AACR Annual Meeting}
    {2025}

  \abstr
    {E.M. Chan, P.C. Borck, N. Bick, K. Jankovic, I. Boyle, L. Parker-Burns, K. Foster, \ldots, \me, \etal}
    {PELO is a synthetic lethal target in chromosome 9p21-deleted or MSI cancers}
    {Cancer Research 84 (6 Suppl.), Abstract 3364; AACR Annual Meeting}
    {2024}

  \abstr
    {J.M. Krill-Burger, \me, B.R. Paolella, J.M. McFarland, F. Vazquez}
    {Systematic methods to identify cancer vulnerabilities from genome-wide loss-of-function screens: an interactive framework for target discovery}
    {Cancer Research 82 (12 Suppl.), Abstract 1897; AACR Annual Meeting}
    {2022}

  \abstr
    {L.D. Cervia, T. Shibue, B. Gaeta, \me, L. Leung, N. Li, N. Dumont, \etal}
    {A ubiquitination cascade regulates the integrated stress response and epithelial cancer survival}
    {Cancer Research 81 (13 Suppl.), Abstract 1950; AACR Annual Meeting}
    {2021}

  \abstr
    {S. Raghavan, P.S. Winter, A.W. Navia, H.L. Williams, A. DenAdel, R.L. Kalekar, \ldots, \me, \etal}
    {Transcriptional subtype-specific microenvironmental crosstalk and tumor cell plasticity in metastatic pancreatic cancer}
    {Cancer Research 80 (22 Suppl.), Abstract PO-058; AACR Virtual Meeting II}
    {2020}

  \abstr
    {P.S. Winter, S. Raghavan, A.W. Navia, H.L. Williams, A. DenAdel, R.L. Kalekar, \ldots, \me, \etal}
    {Subtype-specific microenvironmental crosstalk and tumor cell plasticity in metastatic pancreatic cancer}
    {Cancer Research 80 (21 Suppl.), Abstract PR03; AACR Virtual Meeting II}
    {2020}

  \abstr
    {P.S. Winter, S. Raghavan, A.W. Navia, H. Williams, J. Galvez-Reyes, \ldots, \me, \etal}
    {Matched metastatic pancreatic ductal adenocarcinoma biopsies and organoid models reveal tumor cell transcriptional plasticity and subtype-specific microenvironmental crosstalk}
    {Cancer Research 80 (11 Suppl.), Abstract PR02; AACR Pancreatic Cancer Special Conference}
    {2020}

  \abstr
    {M. Roberts, L. Kageler, K. Bendinelli, S. Bonner, \me, J. Forrester}
    {The role of EGR1 and AP1 in acute myeloid leukemia cell reprogramming toward cell cycle arrest and apoptosis}
    {Cancer Research 80 (16 Suppl.), Abstract 4678; AACR Annual Meeting}
    {2020}

\sectionListEnd


%-------------------------
% Honors & Awards
%-------------------------
\section{Honors \& Awards}
\sectionList
  \resumeItem{\textbf{Spot Award} ($\times$2), Broad Institute --- for going above and beyond \hfill 2020, 2021}
  \resumeItem{\textbf{Phi Beta Kappa} Honor Society \hfill 2019}
  \resumeItem{\textbf{Best Poster Award}, All-College Science Symposium, Dickinson College \hfill 2019}
  \resumeItem{\textbf{Biology Department Summer Research Grant}, Dickinson College \hfill 2018}
  \resumeItem{\textbf{Dana Research Assistantship}, Dickinson College \hfill 2018}
  \resumeItem{\textbf{Pi Mu Epsilon} --- Mathematics National Honor Society \hfill 2018}
  \resumeItem{\textbf{Upsilon Pi Epsilon} --- Computer Science National Honor Society \hfill 2018}
  \resumeItem{\textbf{Richard Howland Memorial Scholarship} --- awarded to one student for excellence in Computer Science \hfill 2018}
  \resumeItem{\textbf{Jane Hill Prize in Computer Science} --- awarded to one first-year for excellence in CS \hfill 2016}
  \resumeItem{\textbf{Torchbearer Award}, Bhumi --- 6 recipients from a field of 8{,}000 volunteers \hfill 2015}
\sectionListEnd

%-------------------------
% Teaching & Mentoring
%-------------------------
\section{Teaching \& Mentoring}
\sectionList

  \cvHeading
    {Instructor --- R Bootcamp for Biomedical Research}{2022, 2023, 2024}
    {University of California, San Francisco}{San Francisco, CA}
    \bullets
      \resumeItem{Brought the R bootcamp curriculum from the Broad Institute to UCSF and taught the course for three consecutive years to postdocs, clinicians, and graduate students.}
      \resumeItem{Adapted the original cancer-data-science curriculum to a broader biomedical audience; trained $250+$ participants across the three offerings.}
    \bulletsEnd

  \cvHeading
    {Course Creator and Teaching Assistant --- Cancer Program R BootCamp}{2019, 2021}
    {Broad Institute of MIT and Harvard}{Cambridge, MA}
    \bullets
      \resumeItem{Designed and implemented the original curriculum to teach postdocs and graduate students the fundamentals of R for cancer data science.}
      \resumeItem{Trained $\sim$100 participants across the two offerings; received a Spot Award for voluntarily developing the course.}
    \bulletsEnd

  \cvHeading
    {Research Mentor --- UCSF PhD Rotation Students}{2023 -- Present}
    {Gilbert Lab, UCSF / Arc Institute}{San Francisco, CA}
    \bullets
      \resumeItem{Mentored two UCSF PhD rotation students --- one from the Bioinformatics (BMI) program and one from the Tetrad (basic biology) program --- through introductory projects in functional genomics.}
      \resumeItem{Provided experimental and computational guidance and supported development of scientific writing, data analysis, and presentation skills.}
    \bulletsEnd

  \cvHeading
    {Teaching Assistant --- Mathematics and Computer Science}{2016 -- 2019}
    {Dickinson College}{Carlisle, PA}
    \bullets
      \resumeItem{Courses: Introduction to Programming I, Introduction to Programming II, Data Structures.}
      \resumeItem{Facilitated lab sections, debugged student code, and graded assignments.}
    \bulletsEnd

\sectionListEnd

%-------------------------
% Service
%-------------------------
\section{Service}
\sectionList

  \cvHeading
    {Ad-hoc Peer Reviewer}{}
    {Nature Biotechnology}{}

  \cvHeading
    {Co-chair --- CodeRATS}{2020 -- 2021}
    {Broad Institute of MIT and Harvard}{Cambridge, MA}
    \bullets
      \resumeItem{Co-led a community for early-career computational researchers.}
      \resumeItem{Held weekly office hours and facilitated periodic workshops on coding skills and tooling.}
    \bulletsEnd

  \cvHeading
    {Volunteering Coordinator (North India)}{Jun 2014 -- Jul 2015}
    {Bhumi}{New Delhi, India}
    \bullets
      \resumeItem{Designed and implemented an after-school mathematics and science curriculum.}
      \resumeItem{Managed $>$100 volunteers and their projects, contributing 2{,}700 hours over a single year.}
    \bulletsEnd

\sectionListEnd

%-------------------------
% Skills
%-------------------------
\section{Skills \& Methods}
\small
\begin{multicols}{2}
\noindent\textbf{Functional genomics:} CRISPR screens (knockout, interference, activation), Perturb-seq, scRNA-seq, viral vector design, mammalian cell culture, cloning.

\vspace{4pt}
\noindent\textbf{Single-cell analysis:} Scanpy, Seurat, Cell~Ranger, scVI, rapids-singlecell (GPU-accelerated workflows), MAGeCK.

\vspace{4pt}
\noindent\textbf{Programming \& ML:} Python, R, PyTorch, JAX. Statistical modeling, dimensionality reduction, foundation models for DNA / RNA / protein.

\columnbreak

\noindent\textbf{AI tooling:} LLM-based agents, multi-agent workflows, retrieval-augmented systems, LangChain, DSPy, scientific APIs, prompt engineering.

\vspace{4pt}
\noindent\textbf{Languages (spoken):} English (native), Hindi (professional), Assamese (conversational).
\end{multicols}
\normalsize

\end{document}
